\section{Discussion}

\subsection{What the ladder does and does not license}

The decomposition reported in \S\ref{sec:res-ladder} is a statement about one class of
models: fixed recurrent reservoirs with a trained linear readout, evaluated on one
spike-based classification task at matched size, density and spectral radius. Within
that class, the attribution is clean, because adjacent rungs differ in exactly one
structural property and nothing else. Outside it, five caveats apply, and only one of them
can be resolved by more instances of the same experiment.

First, the readout is linear and the recurrent weights are untrained. A structural
property might matter for \emph{learnability} -- how readily a network is shaped by a
learning rule -- while being invisible to a fixed-reservoir probe. Motif structure has
been argued to act as an inductive bias, and an inductive bias is by definition a claim
about learning trajectories, not about the function computed by a frozen network. Our
result does not test that claim, and the natural next experiment is to run the same
ladder with the recurrent weights trained by surrogate-gradient descent
\cite{neftci2019surrogate,zenke2021remarkable}, where the rungs would differ in
initialisation rather than in final connectivity. We flag one trap in that experiment,
having walked into it: with only a fraction of neurons driven, a firing-rate penalty
strong enough to be useful in an all-to-all network makes total silence the global
optimum, and a collapsed network returns a perfectly flat ladder that is
indistinguishable from the result reported here. A trained ladder should therefore
report its zero-recurrence control and its mean firing rate before it reports anything
else.

Second, all weights are binary up to valence. Real synapses have graded strengths that
correlate with topological position, and a motif census computed on the binarised graph
discards that correlation. A weighted motif-preserving null is a natural extension; the
annealed chain generalises to weighted targets without modification, but the
degree-preserving move set does not.

Third, the motif-preserving chain is annealed, and annealing is not proven to sample
uniformly over the constraint set. We report mean pairwise overlap between ensemble
members as evidence that the chain explores that set rather than collapsing onto one
region of it (Table~\ref{tab:nullval}), and we show in \S\ref{sec:res-robust} that the
conclusion survives an eightfold variation in chain length and does not depend on which
of the two possible starting points the chain uses. Those checks rule out the two
specific ways an annealed sampler could manufacture a null contrast; they are not a
uniformity guarantee, and a formally exact sampler for a census-constrained
degree-preserving ensemble remains open. We note what would be at stake if the ensemble
were meaningfully non-uniform: the bias would have to correlate with reservoir accuracy
to produce the flat ladder we observe, and the chain's objective makes no reference to
the task, the input projection or the readout.

Fourth, the null contrasts are bounded by power rather than by measurement: at
$\NinstancesLadder$ surrogates per rung the design resolves $d = \MDEHolm$ at the Holm
worst case, so a moderate motif effect would not have been detected. We report the bound
in \S\ref{sec:res-ladder} rather than presenting the null as though it excluded any
effect at all.

Fifth, the benchmark is a 20-class classification with a one-second horizon. Tasks that
demand memory beyond the intrinsic time constants, or that require a specific dynamical
regime, might be more sensitive to recurrent structure. We partly address this with
task-independent dynamical probes (\S\ref{sec:res-dynamics}), which agree with the task
result, but agreement across two measures is not agreement across all of them.

A further objection applies specifically to the degree-sequence result, and we checked
it rather than argued it. Three of the five networks are induced subgraphs grown by
breadth-first expansion from an above-median-degree seed, and growing outward from a hub
is exactly the sampling that could concentrate high-degree nodes and so inflate the
property the effect is attributed to. Measured as the coefficient of variation of total
degree, extraction does not concentrate it: the ratio of subgraph to full-release
heterogeneity is at most $\MaxCVRatio$ across all five networks, and as low as
$\MinCVRatio$ among the $\NExtractedNetworks$ that are genuinely subsampled, with the
remaining $\NWholeNetworks$ used whole. If the extraction biases the degree result at
all, it does so against it.

\subsection{What a unit model decides}

The result of \S\ref{sec:res-unit} is the one we would most like to be wrong about, so we
have tried to state it at exactly its strength and no further.

What it establishes is narrower than the draft of this section that preceded
\S\ref{sec:res-nullchoice}, and we have rewritten it rather than leaving the stronger
phrasing standing. That draft said an advantage unambiguous in graded units was absent in
spiking ones. Against the published null there is no unambiguous advantage in graded units
to be absent: what remains is that the two model classes do not give the same answer on
identical graphs, and that capacity in the spiking model is compressed by a factor of
$\UnitCapRatio$, so effects of the reported size sit near the resolution of the
measurement there. It does not establish that connectome topology is useless in spiking
systems --- one network in our set is at the capacity floor and carries no information
either way --- nor that memory capacity is the right currency.

The mechanism is the one claim in this paper that the null correction moved in our favour,
and it is worth saying so explicitly because we had already published the opposite.
Measured against the ensemble we first built, differential silencing did not predict the
capacity gap and we wrote it up as a failed prediction. Measured against the published
ensemble it does ($r = \UnitSilenceGapR$, $p = \UnitSilenceGapP$, $n =
\UnitSilenceGapN$). We hold the surviving version to the same standard as the discarded
one: it accounts for part of the spiking arm's behaviour in the predicted direction and
leaves most of the variance unexplained.

What we think it does establish is a constraint on how the result should be used. The
connectome-reservoir finding is cited as evidence that biological wiring is computationally
special, and that inference is doing work in proposals for neuromorphic hardware built on
connectome-derived connectivity. Neuromorphic hardware is spiking. Whatever the advantage
measured in $\tanh$ units is worth, the inference from it to a spiking substrate needs its
own evidence --- and on our measurements that evidence does not exist in either direction
yet, because the comparison is weaker in the spiking model rather than resolved there.

We would put the general form of this more plainly than the specific result. A reservoir
result is a statement about a triple: a graph, a unit model and an operating point. Two of
those three are usually chosen once and reported as fixed --- and the null, a fourth, is
usually named rather than specified. We varied all four and the conclusion turns on each.
The methodological recommendation is unglamorous and cheap: report the unit model as a
variable rather than an implementation detail, sweep the gain rather than picking a point
on it, and ship the null-generating code.

\subsection{What the gain axis does to a literature}

The result of \S\ref{sec:res-operating} is uncomfortable in a specific way, and we want to
state the discomfort precisely rather than soften it.

It is not that the published connectome advantage is wrong. It is that the quantity being
reported is a point on a curve whose sign changes, and the curve has not been shown. Our
own attempt to recover the published direction at the published operating point succeeds
against one null ensemble and fails against the other (\S\ref{sec:res-nullchoice}), which
is a reason to doubt our replication rather than the original --- and which makes the
point of this section sharper, not weaker: a single evaluation carries no information
about where on the curve, or against which ensemble, it was taken. Nothing in a
single-operating-point comparison indicates which side of the crossing it falls on, and
nothing in the reported statistics --- ours or anyone's --- carries that information,
because the statistics are computed conditional on the choice.

We think the reason this has gone unremarked is that criticality is a principled choice.
There are good arguments that a reservoir computes best near the edge of chaos
\cite{bertschinger2004real,legenstein2007edge,boedecker2012information}, and if biological
networks operate there, comparing at criticality is comparing where it matters. The
difficulty is that the argument licenses the operating point but not the inference: a
topology can be better at criticality and worse elsewhere without either fact being the
more fundamental one, and ``connectome topology confers computational function'' is a
claim of the second kind stated on evidence of the first.

There is also an asymmetry in the direction of the effect that we think is the more
interesting result. The published comparison is against a \emph{degree-preserving} null.
Against a fully unstructured graph of the same size and density, the connectome does not
lead at any gain we tested. So the advantage that appears at criticality is an advantage
over graphs that share the connectome's degree sequence --- and the degree sequence is,
by our ladder, the single most costly structural property on this task. The narrow win
at criticality sits on top of a large loss that the degree-preserving control holds
constant and therefore hides.

What we would ask of a future comparison of this kind is small: report the contrast across
a range of gains, and state the range over which the sign is stable. That is one sweep
rather than one evaluation, and it converts a claim about a point into a claim about a
regime. We have tried to hold our own results to it.

\subsection{Why matched controls change the conclusion}

The reason a control ladder matters here is arithmetic rather than philosophical. An
Erd\H{o}s--R\'enyi comparison bundles every structural difference into a single contrast,
so whatever it finds gets attributed to whichever property the authors were already
interested in. In our measurements that bundled effect is not small --- the
Erd\H{o}s--R\'enyi rung leads the connectome by up to $\RmLzeroMinusLfour$ of accuracy
in percentage points at
a matched operating point --- which makes the attribution question sharper rather than
moot. Separating the properties shows that the largest single share is spent at the step that
imposes the degree sequence, and that the motif census contributes a further effect of
comparable size --- in the same, negative direction. An Erd\H{o}s--R\'enyi comparison on
this benchmark delivers a large, significant, correctly-signed aggregate effect that a
motif-focused reading would have assigned to motifs. Part of it is indeed the motif
census; most of it is not, and none of it points the way that reading would require. That
is the concrete cost of the bundled control, and it is why we regard the ladder as the
contribution here rather than the result it produces.

This is also where the protocol control in \S\ref{sec:task} does the most work. A
reservoir comparison in which a zeroed recurrent matrix performs as well as a connectome
is not a weak experiment; it is a measurement of the input projection. We report both
protocols precisely because the difference between them dwarfs every structural effect we
measure, and because the failure is silent: nothing in the naive protocol's output
announces that recurrence is unused. We would encourage the same control in any
study comparing connectome-derived reservoirs.

\subsection{Degeneracy is a statement about directions, not magnitudes}

The pruning results in \S\ref{sec:res-degeneracy} make a point that we think is
under-appreciated. Robustness to random ablation is routinely reported as evidence that
biological circuits are degenerate, redundant or gracefully degrading. But a randomly
chosen perturbation direction in a high-dimensional space is nearly orthogonal to any
fixed functional subspace, so tolerance to random damage is close to a generic property
of large recurrent systems rather than a distinguishing property of brains. The
operational alignment $\gamma$ makes part of this quantitative: the same perturbation
cardinality can be almost harmless or markedly damaging depending only on its direction,
and $\gamma$ ranks the directions correctly. It does not deliver the scaling
Proposition~\ref{prop:tolerance} predicts --- the product $\gamma\cdot\|E\|_F^{\max}$
should be constant across directions and instead varies by a factor of
$\GammaProductSpread$ (\S\ref{sec:res-degeneracy}) --- so $\gamma$ should be read as a
diagnostic of which directions are dangerous rather than as a calibrated predictor of how
much damage they do. A degeneracy claim is nonetheless only as strong as the worst
direction it has been tested against, and that conclusion needs only the ordering.

We also note what replacing Davis--Kahan buys. The eigenvector bound is not merely loose
on these networks; it is inapplicable, because the spectra are near-degenerate and the
eigenvectors it bounds are themselves ill-conditioned. Quoting the ratio between an
empirical tolerance and a vacuous bound as a measure of biological robustness -- as an
earlier revision of this work did -- attributes to the system a property of the tool.
The transfer-map bound is weaker in the usual sense (it is first-order and
horizon-limited) but it is about a quantity that task performance actually depends on,
and it requires no gap.

\subsection{A method sentence is not a method}
\label{sec:disc-null}

The single most consequential finding in this paper is also the least glamorous: the
phrase ``randomised while preserving density and degree sequence'', which appears in some
form in nearly every paper in this literature, does not specify an ensemble. Two
constructions satisfying it gave the connectome's advantage opposite significant verdicts
in $\NNullVerdictFlip$ of $\NNullCells$ cells (\S\ref{sec:res-nullchoice}). The ambiguity
is not exotic. It is whether a weight is a property of an edge or a property of the
network's weight distribution, and both readings are natural enough that the phrase has
carried them both for two decades.

We are the demonstration and not the auditor here. Our first implementation was the
rank-reassigning one, written from the published description without reading the toolbox
source, and it produced a result that was clean, large, significant in every network and
in the direction we expected. Nothing about it looked wrong, because nothing about it was
wrong except the question it was answering. Had the two ensembles happened to agree, we
would have reported the number and moved on, and the paper would have contained an
undetectable error of exactly the kind we spent the rest of it cataloguing.

Two consequences follow, and neither is a matter of care. The first is that a replication
which constructs its null from a published description is testing a different hypothesis
from the original, and cannot tell which of the two it has tested. The second is that the
remedy is not more careful prose --- we do not think a longer sentence would have saved
us --- but the null-generating code as an artefact, cited and versioned the way a dataset
is. Where the original code exists, as it does here, calling it is not a courtesy to the
original authors; it is the only way to know what was compared.

This generalises past reservoir computing. Motif enrichment, small-worldness, rich-club
coefficients and every other statistic in network neuroscience defined as a deviation from
a null inherit whatever ambiguity the null's description carries. The $Z$-scores we report
in \S\ref{sec:res-motifs} are no exception, and we state the construction we used in the
methods for that reason.

\subsection{On correcting the record}

A substantial part of this work consisted of removing claims from our own earlier
manuscript: fabricated commercial-process silicon figures, a human microscale tier built
on a connectivity graph that was empty, sign-attributed motif results for datasets with
no usable inhibition, p-values below the resolution of the null ensemble that produced
them, and an ablation table whose own control row contradicted the conclusion drawn from
it. We report these in full in the accompanying audit rather than quietly omitting them,
because the corrected analysis is only interpretable against the record of what it
replaces, and because several of the errors are of kinds that are easy to reproduce:
a silently swallowed exception in a loader, a null model believed to impose a constraint
that it does not, and a benchmark protocol under which the object of study makes no
difference to the result.

The replacement analysis produced four more of the same kind, and we record those too.
The largest is the null-ensemble substitution discussed above, which is not a coding
error at all --- the code did exactly what it was written to do --- but a
misreading of a method sentence that no test could have caught.
Our annealing energy was normalised by a quantity several times larger than the one the
methods described, which left the chain unable to hold its census on the most reciprocal
network in the set. Our ``weakest-first'' pruning arm was executing as whole-neuron
ablation in index order, because these graphs are unweighted and every edge magnitude is
therefore equal. And the first version of Proposition~\ref{prop:gapfree} was stated in
the spectral radius rather than the operator norm, an inequality that holds only for
normal matrices --- in a paper whose central argument is that eigenvalue-based reasoning
fails on these networks. Each was caught by checking a diagnostic that the code already
recorded, or by checking the code against the prose; none was caught by the result
looking wrong, because in every case it did not. We take that to be the substantive
methodological point, and it is why the convergence diagnostics, the protocol control
and the held-out statistic are reported in the paper rather than kept in the
repository.

\subsection{Coupling is a relation, not a property}
\label{sec:disc-coupling}

It is worth stating plainly what the four results have in common, because it is the
thing the title names. Computational coupling --- the degree to which what a network
computes is attributable to how it is wired --- reads like a property of the network, and
is written down that way: one connectome, one number, one verdict. Definition~\ref{def:coupling}
makes explicit that it is nothing of the kind. $\kappa$ takes four arguments besides $W$,
and this paper is an empirical demonstration that none of them is inert. Change
$\mathcal{N}$ and the sign flips in $\NNullVerdictFlip$ of $\NNullCells$ cells with
significance on both sides. Change $U$ and the whole capacity scale compresses by
$\UnitCapRatio\times$, taking the contrast with it. Move along $\rho$ and the sign flips
again, in $\NGainSignFlip$ of $\NGainNetworks$ networks. Choose $R$ badly and the coupling
is identically zero: deleting the recurrent matrix costs nothing
($\ZeroRecFull$ against $\ZeroRecFullConn$), which is to say the measurement was never
touching $W$.

The practical consequence is a reporting standard rather than a new experiment. A
coupling is quotable only with its protocol attached --- and, given
\S\ref{sec:res-nullchoice}, with the null-generating code attached rather than a
description of it, since the description does not determine the ensemble. Where a claim
is meant to be about the connectome rather than about one cell of the design, the
defensible object is not a number but a surface: $\kappa$ over the ranges of
$\mathcal{N}$, $U$, $\rho$ and $R$ that the claim is supposed to survive. That is
why we report axes rather than points throughout, including for our own null results,
and it is why a non-detection here is reported as a non-detection rather than as a
refutation.

We do not claim the four are exhaustive. They are the four our design happened to vary,
and the honest reading of Definition~\ref{def:coupling} is that any further argument
smuggled into $C$ --- the task, the horizon, the delay structure, the regularisation of
the readout --- is a further candidate for the same treatment. The claim we do make is
narrower and, we think, secure: a connectome-versus-null result reported as a single
signed number is under-determined by its own methods section, and on our data the
under-determination is large enough to reverse the conclusion.

\subsection{Outlook}

The motif-preserving ensemble is not specific to neuroscience. Any field that infers
function from subgraph overrepresentation -- gene regulatory networks, food webs,
software dependency graphs -- faces the same attribution problem, and the same annealed
construction applies wherever a motif census admits incremental evaluation. The
immediate extensions we regard as most informative are the trained ladder described above,
applying $\gamma$ to measured biological perturbations --- where the prediction is that
developmental and activity-dependent rewiring occupies the low-$\gamma$ regime that makes
degeneracy cheap --- and extending the null-sensitivity analysis of
\S\ref{sec:res-nullchoice} beyond the two constructions we happened to write. We suspect
the two we compared are not the only pair in circulation under that description, and the
cheapest useful contribution anyone could make to this literature is an inventory of
which published results were measured against which ensemble.
